\documentclass[conference]{IEEEtran}
\IEEEoverridecommandlockouts

% Font: Latin Modern — the LaTeX-native Computer Modern family.
% IEEEtran defaults to Times; these lines override it.
\usepackage[T1]{fontenc}
\usepackage{lmodern}
\renewcommand{\rmdefault}{lmr}
\renewcommand{\sfdefault}{lmss}
\renewcommand{\ttdefault}{lmtt}
\renewcommand{\familydefault}{\rmdefault}

\usepackage{cite}
\usepackage{amsmath,amssymb,amsfonts}
\usepackage{algorithmic}
\usepackage{graphicx}
\usepackage{textcomp}
\usepackage{xcolor}
\usepackage{booktabs}
\usepackage{url}
\usepackage[hidelinks]{hyperref}

\def\BibTeX{{\rm B\kern-.05em{\sc i\kern-.025em b}\kern-.08em
    T\kern-.1667em\lower.7ex\hbox{E}\kern-.125emX}}

\begin{document}

\title{Cost-Sensitive Fraud Decisioning Under Delayed Labels:\\
A Pre-Registered Falsification Study of Operational Robustness}

\author{\IEEEauthorblockN{Ken Ira Lacson Talingting}
\IEEEauthorblockA{\textit{Independent Research}\\
Manila, Philippines \\
\texttt{keniralacson@gmail.com}}}

\maketitle

\begin{abstract}
Fraud detection systems must decide whether to approve, review, or block a
transaction in real time, yet fraud labels arrive days to months later and
are frequently censored before an evaluation window closes. We frame this
as a cost-sensitive decision problem rather than a classification problem
and evaluate it end-to-end on the Bank Account Fraud (BAF) dataset under a
simulated one-month label delay. A Logistic Regression classifier feeding an
argmin-of-expected-cost policy reduces realized cost per transaction by
57.91\% relative to the strongest baseline, LightGBM with a static 0.5
threshold, with a 95\% bootstrap confidence interval of [53.95\%, 62.13\%].
We then subject the result to a pre-registered falsification study across
six hypotheses that test its operational robustness: finite review capacity,
a two-month label delay, rolling retraining, feature drift, group disparity,
and stratified calibration. Three hypotheses survive, one is confirmed with
a fragility note, and two are killed as non-findings; no kill criterion is
modified after results are seen. The headline is not overturned. It is
strengthened by capacity-robustness and delay-robustness, confirmed but
narrowed under rolling retraining, and bounded below materiality by
stratified calibration. We report censored-label counts ($96{,}843$ of
$1{,}000{,}000$, $9.68\%$) rather than treating unobserved fraud as
legitimate, and we release the full evaluation protocol and falsification
plan alongside the pipeline. The paper contributes a reproducible template
for evaluating fraud decision policies under delayed and partial labels, and
demonstrates that static thresholds on highly imbalanced data behave
near-identically to approve-all even when the underlying model is
well-calibrated.
\end{abstract}

\begin{IEEEkeywords}
cost-sensitive learning, fraud detection, delayed feedback, censored labels,
decision policy, temporal backtest, pre-registered falsification,
operational robustness, Bank Account Fraud dataset, expected-cost
minimization, bootstrap confidence intervals, sensitivity analysis
\end{IEEEkeywords}

% =============================================================
\section{Introduction}
% =============================================================

Digital payment systems process millions of transactions every day, and the
operational question a fraud system faces is not whether a transaction is
fraudulent in hindsight, but whether to approve, review, or block it
\emph{before} the label arrives. This asymmetry is structural. Authorization
happens in milliseconds; disputes and chargebacks resolve days to weeks
later, and some fraud is never reported at all. A supervised model trained
on settled disputes therefore learns from a label that arrives late and is
censored not-at-random.

Most published work in fraud detection treats the problem as binary
classification and reports ranking metrics such as area under the ROC curve
(AUC) or area under the precision--recall curve (PR-AUC). These metrics
compare models on a fixed labeled dataset. They do not answer the
operational question: given a score, should this transaction be approved,
reviewed, or blocked, and what is the realized cost of that decision? On a
dataset where fraud is roughly $1\%$ of transactions, a high AUC can coexist
with a policy that barely improves on approving everything.

This paper takes a different framing. We treat fraud detection as a
one-step decision problem under delayed, partial, and asymmetric-cost
feedback. The model produces a probability, but the system under evaluation
is the \emph{policy} that maps that probability, the transaction amount, and
a frozen cost matrix into one of three actions. The policy chooses the
action that minimizes expected cost. This is not a new policy ---
expected-cost minimization is standard in cost-sensitive learning --- but it
is rarely evaluated end-to-end on a fraud dataset with simulated label
delay, censored labels reported rather than hidden, and a pre-registered
falsification study on its operational robustness.

We evaluate this framing on the Bank Account Fraud (BAF) dataset, a public
one-million-row synthetic benchmark designed for fraud research. Because BAF
provides only month-level time granularity, we simulate a one-month label
delay and hold the final month out as censored. We compare a cost-sensitive
policy against six baselines under an identical temporal split and cost
matrix, report realized cost per transaction as the primary metric, and test
robustness with a $2\times$ sensitivity sweep across cost parameters and a
bootstrap confidence interval.

We then test whether the headline survives operational reality. Six
falsifiable hypotheses are pre-registered with numeric kill criteria written
before the tests run. The hypotheses cover review capacity, multi-regime
delay, rolling retraining, feature drift, group disparity, and stratified
calibration. Two of the six are killed and reported as non-findings. The
headline is not overturned by any test.

\subsection{Contributions}

\begin{itemize}
  \item A reproducible evaluation protocol for fraud decision policies under
  delayed and censored labels, including pre-registered stop criteria and an
  explicit prohibition on random splits.
  \item A temporal backtest on BAF with six canonical baselines,
  censored-label reporting, cost sensitivity, and bootstrap confidence
  intervals. The policy reduces realized cost per transaction by $57.91\%$
  relative to the strongest baseline, with a $95\%$ CI of $[53.95\%,
  62.13\%]$.
  \item A pre-registered falsification study across six operational
  hypotheses. Three survive, one is confirmed with a fragility note, two are
  killed. The headline is strengthened by capacity-robustness and
  delay-robustness, narrowed under rolling retraining, and bounded below
  materiality by stratified calibration.
  \item A documented negative result at the framing level: static thresholds
  on highly imbalanced fraud data are not decision rules. On this dataset,
  the strongest static threshold reduces cost by only $5.97\%$ relative to
  approve-all.
\end{itemize}

\subsection{Organization}

Section~\ref{sec:related} reviews related work. Section~\ref{sec:problem}
states the problem from first principles. Section~\ref{sec:policy} defines
the decision policy. Section~\ref{sec:data} describes the dataset,
delay simulation, temporal split, and cost matrix.
Section~\ref{sec:setup} describes the experimental setup and baselines.
Section~\ref{sec:results} reports the temporal backtest.
Section~\ref{sec:robust} reports the falsification study.
Section~\ref{sec:discussion} discusses implications. Section~\ref{sec:lim}
states limitations. Section~\ref{sec:conclusion} concludes.

% =============================================================
\section{Related Work}
\label{sec:related}
% =============================================================

\subsection{Cost-Sensitive Learning}

Elkan \cite{elkan2001} formalized the theory of cost-sensitive learning and
showed that a classifier trained to minimize misclassification error is
suboptimal when class costs are asymmetric. The decision-theoretic view ---
choose the action that minimizes expected cost --- is the foundation of our
policy. Our work differs in the setting: we apply the policy to a fraud
dataset under simulated label delay, and we evaluate it end-to-end rather
than comparing classifiers on a fixed labeled split.

\subsection{Delayed Feedback}

Chapelle \cite{chapelle2014} studied delayed feedback in display
advertising, where conversions arrive long after impressions. That work
introduces the delay-aware evaluation problem and provides a probabilistic
model for the conversion-time distribution. We adopt the delay-aware
evaluation structure but do not model the delay distribution: BAF exposes
month-level granularity only, so we fix a one-month regime and hold the
final month out as censored. A two-month regime is tested in
Section~\ref{sec:h2}.

\subsection{Fraud Detection Datasets}

The Bank Account Fraud (BAF) suite \cite{jesus2022} is a NeurIPS 2022
benchmark designed to be biased, imbalanced, and dynamic, providing
month-level temporal information that supports delay-aware evaluation. It is
synthetic, which we treat as a limitation rather than a feature; its value
here is that it is large enough to train reliably, is licensed for
redistribution, and exposes a time column that is required for temporal
splits.

\subsection{Calibration Under Imbalance}

Dal Pozzolo et al.\ \cite{dalpozzolo2015} demonstrated that
class-imbalanced fraud classifiers often produce uncalibrated scores and
that calibration materially affects downstream decisions. Our policy
assumes $p = P(\text{fraud} \mid X_t)$ is a real probability, so we require
an expected calibration error (ECE) below $0.05$ before a classifier is
admissible as a policy input. In our case the selected classifier passes the
gate without any post-hoc calibration.

\subsection{Positioning}

Our contribution is not a new model or a new policy. It is a demonstration
that a standard cost-sensitive policy, evaluated under a documented temporal
backtest, produces a large and stable cost reduction on an imbalanced fraud
dataset --- and that the reduction survives a pre-registered falsification
study that tests six operational assumptions in sequence. The falsification
discipline, including the reporting of two killed hypotheses at the same
length as the survivors, is the main methodological contribution.

% =============================================================
\section{Problem Formulation}
\label{sec:problem}
% =============================================================

\subsection{The Structural Fact}

In real-time payment systems, authorization happens before the outcome is
known:
\begin{equation}
\text{decision\_time} = t, \qquad \text{label\_time} = t + \Delta_t,
\qquad \Delta_t > 0 .
\end{equation}
This is not a modeling choice. It is the shape of the domain. It cannot be
optimized away by a better classifier, a larger dataset, or a more complex
architecture.

\subsection{Consequences}

Because $\Delta_t > 0$, the label is late, noisy, and sometimes missing. Six
consequences follow: label delay, censoring (some fraud is never reported),
label noise (some disputes are not fraud), class imbalance (fraud is rare),
cost asymmetry (a missed fraud and a blocked customer do not cost the same),
and concept drift (fraud patterns change over time). Each is named in
Section~\ref{sec:lim} as a limitation or tested in
Section~\ref{sec:robust}.

\subsection{The Objective}

For each transaction the system selects the action that minimizes expected
operational cost:
\begin{equation}
a^* = \arg\min_{a \in \{\text{approve, review, block}\}}
\mathbb{E}[\text{cost}(a)] .
\end{equation}
Classification metrics are supporting evidence, not the objective. A
classifier with a higher macro F1 but a worse policy cost is a worse
classifier for this project.

% =============================================================
\section{Decision Policy}
\label{sec:policy}
% =============================================================

The policy consumes a calibrated probability $p$, the transaction amount
$a$, and a frozen cost matrix, and outputs one action. The three expected
costs are:
\begin{align}
\mathbb{E}[\text{cost(approve)}] &= p \cdot L_{\text{fraud}}(a), \label{eq:app}\\
\mathbb{E}[\text{cost(review)}]  &= c_{\text{rev}} + p \cdot r_{\text{fraud}}, \label{eq:rev}\\
\mathbb{E}[\text{cost(block)}]   &= (1 - p) \cdot c_{\text{fp}}, \label{eq:blk}
\end{align}
where $L_{\text{fraud}}(a)$ is the fraud loss if approved, $c_{\text{rev}}$
is the review cost paid regardless of outcome, $r_{\text{fraud}}$ is the
residual fraud loss after review, and $c_{\text{fp}}$ is the false-positive
cost if a legitimate transaction is blocked. When amount scaling is enabled,
$L_{\text{fraud}}(a) = a \cdot r_{\text{loss}}$, where $r_{\text{loss}} =
1 / \overline{a}_{\text{train}}$ is calibrated on training-window amounts
only.

The policy chooses $a^* = \arg\min_{a} \mathbb{E}[\text{cost}(a)]$. The
argmin rule is the source of truth. Derived thresholds on $p$ --- for
example, $\partial \mathbb{E}[\text{cost(review)}] / \partial p = \partial
\mathbb{E}[\text{cost(approve)}] / \partial p$ implies $p_{\text{rev}} =
c_{\text{rev}} / (L_{\text{fraud}} - r_{\text{fraud}})$ --- are a
diagnostic view for reporting, not the implementation.

\subsection{Cost Matrix}

The frozen cost matrix used in all experiments is
\begin{equation}
c_{\text{rev}} = 0.02,\; r_{\text{fraud}} = 0.30,\; c_{\text{fp}} = 0.10,\;
r_{\text{loss}} = 0.0019187869 .
\end{equation}
The rate $r_{\text{loss}}$ satisfies $\overline{r_{\text{loss}} \cdot a} =
1$ on the training window, matching the scale of the constant-loss variant
without using any test-window information.

\subsection{Calibration Requirement}

The policy assumes $p$ is a probability. A classifier with $\text{ECE} \geq
0.05$ on the validation window is not admissible as a policy input unless a
documented calibration step brings it below. Our selected classifier
achieves $\text{ECE} = 0.0033$ on validation and no calibration step is
applied.

% =============================================================
\section{Data and Evaluation Protocol}
\label{sec:data}
% =============================================================

\subsection{Dataset}

The Bank Account Fraud (BAF) suite \cite{jesus2022} is a public
$1{,}000{,}000$-row synthetic benchmark. The positive class rate is
$1.1029\%$ overall. The dataset exposes a month column with values $0$ to
$7$. We use the target \texttt{fraud\_bool} and 28 classifier features,
excluding the row identifier, the time column, the target, the derived delay
columns, and any field that may encode post-decision information. The BAF
license (CC BY 4.0) permits redistribution.

\subsection{Delay Simulation}

BAF provides month-level granularity only, so we fix a one-month delay
regime: $\text{label\_month} = \text{month} + 1$, and a transaction is
observed iff $\text{label\_month} \leq 7$. Transactions in month 7 have
$\text{label\_month} = 8$, beyond the dataset. They are \emph{censored},
never treated as legitimate.

\subsection{Temporal Split}

Table~\ref{tab:split} shows the single chronological split used for the
backtest. No shuffling occurs at any stage.

\begin{table}[!t]
\caption{Chronological split for the one-month delay regime.}
\label{tab:split}
\centering
\small
\begin{tabular}{lrrc}
\toprule
Split & Months & Rows & Fraud rate \\
\midrule
Train      & 0--2 & 397{,}039 & $\approx 1.10\%$ \\
Validation & 3--4 & 278{,}627 & $1.0207\%$ \\
Test       & 5--6 & 227{,}491 & $1.2576\%$ \\
Censored   & 7    & 96{,}843  & --- \\
\midrule
Observed   &      & 903{,}157 & \\
\bottomrule
\end{tabular}
\end{table}

An assertion in the pipeline verifies $\text{train} + \text{validation} +
\text{test} = \text{observed}$, guarding against off-by-one month
boundaries. The test window is used exactly once, after model and policy
are frozen. The test-window fraud rate ($1.2576\%$) is slightly higher than
the training rate ($\approx 1.10\%$). This drift is documented, not
corrected; correcting it would amount to tuning on the test window.

\subsection{Censored Labels}

$96{,}843$ of $1{,}000{,}000$ rows ($9.68\%$) have
$\text{label\_month} > 7$ and are excluded from training and evaluation.
Their count is reported in every results table. Treating them as negatives
would underestimate the fraud rate; treating them as positives would
overstate it.

\subsection{Calibration Gate}

Every classifier is scored on the validation window and its ECE is computed
over ten quantile bins. A classifier is admissible as a policy input iff
$\text{ECE} < 0.05$, or a single documented calibration step brings it
below. Our selected classifier achieves $\text{ECE} = 0.0033$; no
calibration step is applied. Brier on the test window is $0.011506$.

% =============================================================
\section{Experimental Setup}
\label{sec:setup}
% =============================================================

\subsection{Classifiers}

Three traditional classifiers are compared as policy inputs: Logistic
Regression (one-hot encoding plus standard scaling), Random Forest (ordinal
encoding, no scaling), and LightGBM (native categorical handling, no
scaling). Neural networks, deep learning, transformers, and AutoML are
excluded by scope. Class reweighting and resampling are rejected because
they distort calibrated probabilities, and the policy depends on calibration.

\subsection{Hyperparameter Grids}

Each algorithm uses a small documented grid. Logistic Regression varies $C
\in \{0.01, 0.1, 1.0, 10.0\}$ with $\text{max\_iter} = 1000$; Random Forest
varies $\text{n\_estimators} \in \{100, 200, 400\}$; LightGBM varies
$\text{n\_estimators} \in \{100, 300, 500\}$ and $\text{learning\_rate} \in
\{0.05, 0.10\}$. Each configuration is trained under expanding-window
cross-validation by month on the training window; the training window
yields two folds. Selection is by mean realized cost after the policy, not
by macro F1.

\subsection{Noise-Band Guard}

A classifier selection counts as a finding only if the top classifier wins
every fold \emph{and} the relative gap to the runner-up exceeds $5\%$.
Otherwise the result is a non-finding and the tie is broken by simplicity,
in the order Logistic Regression, Random Forest, LightGBM. This guard is
applied before the test window is touched.

\subsection{Baselines}

Six canonical baselines are evaluated under the same split, cost matrix,
and delay regime: (1) random decision (seeded), (2) approve-all, (3)
block-all, (4) Logistic Regression $+$ static $0.5$, (5) Random Forest $+$
static $0.5$, (6) LightGBM $+$ static $0.5$. The policy is compared against
the strongest of the six, not the weakest.

\subsection{Statistical Rigor}

Every cost comparison is accompanied by a $95\%$ bootstrap confidence
interval on $1{,}000$ resamples with seed $42$. An improvement counts as a
finding only if the interval on the advantage excludes zero \emph{and} the
relative cost reduction exceeds $5\%$.

% =============================================================
\section{Results: Temporal Backtest}
\label{sec:results}
% =============================================================

\subsection{Classifier Selection}

Table~\ref{tab:cv} shows cross-validated realized cost per transaction for
the three classifiers. The top gap is $1.07\%$, below the pre-registered
$5\%$ noise band, so the comparison is reported as a non-finding and
Logistic Regression is selected by simplicity.

\begin{table}[!t]
\caption{Cross-validated realized cost per transaction on the training
window. Lower is better.}
\label{tab:cv}
\centering
\small
\begin{tabular}{lcc}
\toprule
Classifier & Mean $\pm$ SD & Relative gap \\
\midrule
LightGBM            & $0.005852 \pm 0.000300$ & --- \\
Logistic Regression & $0.005915 \pm 0.000338$ & $1.07\%$ \\
Random Forest       & $0.006515 \pm 0.000579$ & $11.3\%$ \\
\bottomrule
\end{tabular}
\end{table}

\subsection{Policy Versus Baselines}

Table~\ref{tab:backtest} reports realized cost per transaction on the test
window. The policy achieves $0.007491$ cost per transaction; the strongest
baseline is LightGBM with a static $0.5$ threshold at $0.017798$. The
advantage is $57.91\%$ with a $95\%$ bootstrap CI of $[53.95\%, 62.13\%]$,
which excludes zero.

\begin{table}[!t]
\caption{Realized cost per transaction on the frozen test window
($227{,}491$ rows). The policy is the argmin of expected cost.}
\label{tab:backtest}
\centering
\small
\begin{tabular}{lr}
\toprule
Scenario & Cost / txn \\
\midrule
Block-all                     & $0.098742$ \\
Random                        & $0.047484$ \\
Approve-all                   & $0.018928$ \\
Random Forest $+$ static $0.5$& $0.018928$ \\
Logistic Regression $+$ 0.5   & $0.018737$ \\
LightGBM $+$ static $0.5$     & $0.017798$ \\
\midrule
\textbf{Cost-sensitive policy} & $\mathbf{0.007491}$ \\
\bottomrule
\end{tabular}
\end{table}

The static-threshold baselines are informative. The strongest of the three
(LightGBM $+$ $0.5$) reduces cost by only $5.97\%$ relative to approve-all.
Logistic Regression $+$ $0.5$ reduces by $1.01\%$; Random Forest $+$ $0.5$
reduces by $0\%$. On highly imbalanced fraud data, a static threshold is
not a decision rule.

\subsection{Action Distribution}

The policy routes $205{,}395$ transactions to approve, $21{,}371$ to review
($9.39\%$), and $725$ to block. The review band is where most of the cost
reduction lives; the block band is small and precise.

\subsection{Cost Sensitivity}

Each of the three cost parameters is varied over a $2\times$ range.
The policy remains the lowest-cost action at every variation. The minimum
advantage across the sweep is $47.76\%$, at $\text{review\_cost} = 0.04$.
Table~\ref{tab:sens} summarizes.

\begin{table}[!t]
\caption{Cost sensitivity sweep. Each parameter is varied one at a time.}
\label{tab:sens}
\centering
\small
\begin{tabular}{lcc}
\toprule
Variation & Policy cost & Advantage \\
\midrule
Baseline                          & $0.007491$ & $57.91\%$ \\
$c_{\text{fp}} = 0.05$            & $0.007037$ & $60.36\%$ \\
$c_{\text{fp}} = 0.20$            & $0.007625$ & $57.38\%$ \\
$c_{\text{rev}} = 0.01$           & $0.006253$ & $64.87\%$ \\
$c_{\text{rev}} = 0.04$           & $0.009298$ & $47.76\%$ \\
$r_{\text{fraud}} = 0.15$         & $0.006588$ & $62.99\%$ \\
$r_{\text{fraud}} = 0.60$         & $0.008700$ & $51.12\%$ \\
\bottomrule
\end{tabular}
\end{table}

\subsection{Bootstrap Confidence Intervals}

The $95\%$ bootstrap CIs are $[0.007196, 0.007797]$ for the policy cost,
$[0.016967, 0.018670]$ for the LightGBM $+$ $0.5$ baseline, and
$[0.009602, 0.011058]$ for the difference. The advantage CI is
$[53.95\%, 62.13\%]$.

% =============================================================
\section{Operational Robustness: A Pre-Registered Falsification Study}
\label{sec:robust}
% =============================================================

The temporal backtest produces a single, static result under one cost
matrix, one delay regime, and one training split. The result is only as
meaningful as the assumptions that produced it. This section tests whether
the headline survives operational reality: finite review capacity, a longer
label delay, rolling retraining, input drift, group disparity, and
confounding.

\subsection{Method}

Six falsifiable hypotheses are pre-registered with numeric kill criteria
written before the tests run. No kill criterion is modified after results
are seen. A hypothesis that fails its criterion is reported as a
non-finding and dropped from scope. The execution order is cheapest-first,
so that early kills reduce the scope of later work. Non-findings are
reported at the same length and rigor as findings.

\subsection{H1: Finite Review Capacity}

\textbf{Hypothesis.} The policy advantage ($>5\%$ relative cost reduction)
is preserved when review is constrained to a finite per-window capacity.

\textbf{Kill criterion.} At any tested capacity $K \geq 4{,}550$ ($2\%$ of
test volume), the advantage drops below $5\%$.

\textbf{Test.} Re-rank the frozen test-window scores by
$\min(c_{\text{approve}}, c_{\text{block}}) - c_{\text{review}}$, keep the
top $K$ in review, and route the rest by approve-vs-block argmin. No
classifier is retrained; no test score is recomputed.

\textbf{Result.} Table~\ref{tab:h1} reports the advantage at five capacity
levels. The minimum in the kill zone is $44.25\%$ at $K = 4{,}550$; the
margin above the kill threshold is $39.25$ percentage points. Above
$\approx 9.4\%$ capacity the unconstrained advantage ($57.91\%$) is
recovered in full. The policy degrades gracefully as review tightens; there
is no discontinuity.

\textbf{Verdict.} \textbf{Survived.} V1's unconstrained-review assumption is
not load-bearing for the headline.

\begin{table}[!t]
\caption{H1 — policy advantage under constrained review capacity.}
\label{tab:h1}
\centering
\small
\begin{tabular}{rrcc}
\toprule
$K$ & Fraction & Advantage & Kill zone \\
\midrule
$2{,}275$  & $1\%$  & $41.88\%$ & no \\
$4{,}550$  & $2\%$  & $44.25\%$ & yes (min) \\
$11{,}375$ & $5\%$  & $53.11\%$ & yes \\
$22{,}750$ & $10\%$ & $57.91\%$ & yes \\
$45{,}499$ & $20\%$ & $57.91\%$ & yes \\
\bottomrule
\end{tabular}
\end{table}

\subsection{H2: Multi-Regime Delay}

\textbf{Hypothesis.} The advantage is preserved when label delay is
two months instead of one.

\textbf{Kill criterion.} Advantage drops below $5\%$ at a two-month delay.

\textbf{Test.} Re-run the pipeline with $\text{label\_offset} = 2$. The
split becomes train 0--1, validation 2--3, test 4--5, censored 6--7. The
classifier is retrained on the smaller window; the policy and cost matrix
are unchanged.

\textbf{Result.} Censoring rises from $9.68\%$ to $20.50\%$ and the training
window shrinks from $397{,}039$ to $260{,}060$ rows. The policy advantage is
$58.08\%$ with a $95\%$ CI of $[56.35\%, 59.72\%]$, marginally wider than
V1's $57.91\%$ and well above the kill threshold. Validation ECE is
$0.0035$, below the $0.05$ gate.

\textbf{Verdict.} \textbf{Survived.} The decision-centric framing is robust
to the delay regime tested. V1's single one-month regime is not
load-bearing.

\begin{table}[!t]
\caption{H2 — headline comparison across delay regimes.}
\label{tab:h2}
\centering
\small
\begin{tabular}{lcc}
\toprule
Quantity & V1 (1-month) & H2 (2-month) \\
\midrule
Censoring rate           & $9.68\%$             & $20.50\%$ \\
Train rows               & $397{,}039$          & $260{,}060$ \\
Test rows                & $227{,}491$          & $247{,}014$ \\
Policy cost / txn        & $0.007491$           & $0.007289$ \\
Advantage                & $57.91\%$            & $58.08\%$ \\
Bootstrap $95\%$ CI      & $[53.95\%, 62.13\%]$ & $[56.35\%, 59.72\%]$ \\
ECE (validation)         & $0.0033$             & $0.0035$ \\
\bottomrule
\end{tabular}
\end{table}

\subsection{H3: Rolling Retraining}

\textbf{Hypothesis.} Logistic Regression remains the selected classifier
under rolling monthly retraining, as it did under V1's static split.

\textbf{Kill criterion.} The model choice flips in either rolling step.

\textbf{Test.} Two steps are run: $m=5$ (train 0--3, validate 4) and $m=6$
(train 0--4, validate 5). At each step, LR, RF, and LGBM are retrained with
the same grids, and the noise-band guard is applied.

\textbf{Result.} LR is selected at both steps. But the LGBM-vs-LR gap under
rolling retraining ($4.25\%$ and $4.17\%$) is roughly four times larger
than V1's static split ($1.07\%$), and both rolling steps sit within
$0.75$ percentage points of the $5\%$ flip threshold.

\textbf{Verdict.} \textbf{Confirmed, fragile.} The selection stands, but
the margin is thinner than the single static split suggested. This is
recorded as a fragility note rather than a strong confirmation.

\begin{table}[!t]
\caption{H3 — noise-band gap, rolling versus static.}
\label{tab:h3}
\centering
\small
\begin{tabular}{lc}
\toprule
Split & LGBM-vs-LR gap \\
\midrule
V1 static split        & $1.07\%$ \\
H3 step 1 ($m=5$)      & $4.25\%$ \\
H3 step 2 ($m=6$)      & $4.17\%$ \\
Noise-band threshold   & $5.00\%$ \\
\bottomrule
\end{tabular}
\end{table}

\subsection{H4: Feature Drift}

\textbf{Hypothesis.} Feature drift is measurable month-over-month: at least
one feature shows PSI $\geq 0.1$ between consecutive months.

\textbf{Kill criterion.} Maximum PSI across all features and all adjacent
month pairs is below $0.1$.

\textbf{Test.} Compute PSI for all 28 classifier features between adjacent
months on the raw feature table. Binning for continuous features uses
deciles of the reference month. No new training occurs.

\textbf{Result.} Max PSI is $3.9429$ on \texttt{velocity\_4w} between months
0 and 1, far above the $0.1$ kill threshold. 23 of 28 features cross
$\text{PSI} \geq 0.1$ in at least one adjacent pair. Cross-referencing
drift against the classifier's LR importances shows drift is concentrated
on low-importance features: the drifted features together carry $18.84\%$
of importance, only one of the top-5 importance features
(\texttt{payment\_type}, PSI $0.1730$) drifted, and the largest-magnitude
drift (\texttt{velocity\_4w}) carries $0.49\%$ importance.

\textbf{Verdict.} \textbf{Survived.} The static split is challenged by
drift but not invalidated. A monitoring component has signal to detect.

\begin{table}[!t]
\caption{H4 — maximum PSI per adjacent month pair.}
\label{tab:h4}
\centering
\small
\begin{tabular}{rrc}
\toprule
Pair & Max PSI & Feature \\
\midrule
$0 \to 1$ & $3.9429$ & \texttt{velocity\_4w} \\
$1 \to 2$ & $0.5107$ & \texttt{velocity\_4w} \\
$2 \to 3$ & $3.0454$ & \texttt{velocity\_4w} \\
$3 \to 4$ & $0.9611$ & \texttt{velocity\_4w} \\
$4 \to 5$ & $1.6714$ & \texttt{velocity\_4w} \\
$5 \to 6$ & $0.5549$ & \texttt{velocity\_4w} \\
$6 \to 7$ & $3.3890$ & \texttt{velocity\_4w} \\
\bottomrule
\end{tabular}
\end{table}

\subsection{H5: Group Disparity}

\textbf{Hypothesis.} Block rate differs by more than $5$ percentage points
across at least one grouping of transactions by a candidate protected
attribute.

\textbf{Kill criterion.} Maximum block-rate disparity is below $5$
percentage points across all tested groupings.

\textbf{Test.} Compute block rate, review rate, and fraud rate per group for
three groupings: \texttt{customer\_age} in pre-registered bins,
\texttt{employment\_status}, and \texttt{housing\_status}. Groups with $n <
500$ are excluded from disparity calculation.

\textbf{Result.} Table~\ref{tab:h5} reports the disparities. The maximum
block-rate disparity is $2.26$ percentage points on
\texttt{customer\_age}. Review-rate disparities are materially larger
($15.46$, $11.73$, and $33.18$ percentage points across the three
groupings), but they are not the pre-registered metric and are reported as
supporting evidence only.

\textbf{Verdict.} \textbf{Killed.} No block-rate disparity above the
pre-registered threshold is observed. Fairness constraints are not
justified by this diagnostic. The review-rate observation is recorded as a
pre-registration limitation for future work.

\begin{table}[!t]
\caption{H5 — group disparity by grouping (percentage points).}
\label{tab:h5}
\centering
\small
\begin{tabular}{lcc}
\toprule
Grouping & Block-rate & Review-rate \\
\midrule
\texttt{customer\_age}     & $2.26$  & $15.46$ \\
\texttt{employment\_status}& $1.36$  & $11.73$ \\
\texttt{housing\_status}   & $1.76$  & $33.18$ \\
\bottomrule
\end{tabular}
\end{table}

\subsection{H6: Stratified Calibration}

\textbf{Hypothesis.} The direct cost estimate is not materially confounded
by the policy's deterministic selection on classifier inputs.

\textbf{Kill criterion.} Any action stratum has $\text{ECE} \geq 0.05$.

\textbf{Test.} Split the test window by action (approve, review, block) and
compute ECE within each stratum. Under conditional exchangeability given
the model's inputs, per-stratum calibration is a necessary condition for
the direct estimate to be unbiased.

\textbf{Result.} The block stratum shows $\text{ECE} = 0.0600$, above the
$0.05$ threshold. The stratum is $725$ rows ($0.32\%$ of the test window)
and carries $24.02\times$ the aggregate ECE. The materiality check
required by the H6 kill clause bounds the plausible bias on the headline
advantage below the $5\%$ threshold that would authorize causal machinery.

\textbf{Verdict.} \textbf{Killed (immaterial).} The block stratum is
miscalibrated, but the stratum is too small for the miscalibration to
materially affect the headline. No IPW, doubly robust, or randomized
logging work is authorized.

\begin{table}[!t]
\caption{H6 — per-stratum calibration on the test window.}
\label{tab:h6}
\centering
\small
\begin{tabular}{lrrrr}
\toprule
Stratum & $n$ & Base & $\overline{p}$ & ECE \\
\midrule
approve & $205{,}395$ & $0.69\%$ & $0.0059$ & $0.0011$ \\
review  & $21{,}371$  & $5.65\%$ & $0.0422$ & $0.0144$ \\
block   & $725$       & $32.69\%$& $0.2794$ & $\mathbf{0.0600}$ \\
\bottomrule
\end{tabular}
\end{table}

\subsection{Summary of Verdicts}

\begin{table}[!t]
\caption{Pre-registered falsification study: six verdicts.}
\label{tab:verdicts}
\centering
\small
\begin{tabular}{lll}
\toprule
ID & Hypothesis & Verdict \\
\midrule
H1 & Capacity-constrained        & survived \\
H2 & Two-month delay             & survived \\
H3 & Rolling retraining          & confirmed (fragile) \\
H4 & Feature drift               & survived \\
H5 & Group disparity             & killed \\
H6 & Stratified calibration      & killed (immaterial) \\
\bottomrule
\end{tabular}
\end{table}

The V1 headline is not overturned by any test. It is strengthened by H1 and
H2, confirmed with a fragility note by H3, challenged but not broken by H4,
and bounded below materiality by H6. Two proposed concerns (H5 and H6) are
tested and dismissed.

% =============================================================
\section{Discussion}
\label{sec:discussion}
% =============================================================

\subsection{Why the Policy Wins}

The static-threshold baselines are informative. LightGBM with a static
$0.5$ cutoff reduces realized cost by only $5.97\%$ relative to
approve-all. Logistic Regression $+$ $0.5$ reduces by $1.01\%$, and Random
Forest $+$ $0.5$ reduces by $0\%$. On this dataset, a well-calibrated
classifier with a fixed threshold is nearly indistinguishable from
approving every transaction. The policy's advantage comes from the cost
rule, not from the classifier. This is the paper's most operationally
useful finding.

\subsection{What the Falsification Study Shows}

The falsification study is the paper's main methodological contribution.
Each of V1's load-bearing assumptions is tested in isolation with a numeric
kill criterion written before the test. Three assumptions are confirmed as
robust, one is narrowed (classifier selection is confirmed but fragile under
rolling retraining), and two proposed concerns are dismissed. Every verdict
is reported at the same length as the others. Two of the six are
non-findings.

The falsification discipline is what makes the robustness claim credible.
Without it, "the result survives X" is an unfalsifiable statement. With it,
the reader can see exactly which conditions were tested, what threshold
would have killed each, and what the measured value was.

\subsection{Capacity, Delay, and Drift}

H1 shows the policy advantage survives a $4.7\times$ reduction in review
capacity below the observed band size. H2 shows the advantage survives a
doubling of the label delay. H4 shows drift is real but concentrated on
features the classifier does not rely on. Together these three results
suggest that the decision-centric framing is not an artifact of a favorable
split, favorable delay, or favorable feature stability.

\subsection{The Rolling-Retraining Fragility}

H3 is the one place where the study finds a narrowing. The classifier
selection is stable under rolling retraining, but the gap between the top
two classifiers is four times larger than under the static split. The
selection stands; the margin is thinner than V1 alone implied. This is a
caveat, not a reversal, and it is reported as such.

\subsection{Two Non-Findings}

H5 and H6 are killed. Neither fairness constraints nor causal machinery are
justified by the diagnostics. Reporting these as non-findings, at the same
length as the survivors, is deliberate. A falsification study that reports
only the survivors is not a falsification study.

% =============================================================
\section{Limitations}
\label{sec:lim}
% =============================================================

\subsection{Dataset Limitations}

BAF is synthetic. Real fraud distributions differ in ways that the
benchmark cannot capture. The reported cost reduction is not a production
estimate. Fraud is assumed to occur at a fixed overall rate of $1.1029\%$;
real systems face drifting rates.

\subsection{Delay Regime}

Only one- and two-month delay regimes are tested. The dataset provides
month-level granularity only; sub-month regimes are not possible. A
production system faces delay distributions, not point delays.

\subsection{Amount Proxy}

The amount column is proxied by \texttt{proposed\_credit\_limit}. It
correlates with transaction size but is not a true transaction amount, and
it does not scale false-positive cost.

\subsection{Capacity Model}

Capacity is modeled as a per-window count of reviews. Wall-clock latency,
investigator skill variation, and queue dynamics are not modeled.

\subsection{Feedback and Adversarial Effects}

BAF is a static dataset. The policy's decisions do not affect future
labels. Fraudsters are assumed not to adapt. Both effects are real in
production and are named here rather than modeled.

\subsection{Pre-Registration Limitations}

The H5 kill criterion used absolute block-rate disparity, which is bounded
below $1$ percentage point by the policy's overall block rate. Review-rate
disparity is materially larger. The H6 kill criterion applies at the
stratum level without a sample-size correction; the block stratum is $725$
rows and ECE over such a stratum has substantial sampling noise. Both
criteria are applied as written.

\subsection{What Is Not Claimed}

Production readiness is not claimed. Generalization to other fraud domains
is not claimed. A higher AUC than prior work is not claimed. Causality is
not claimed. The results are specific to the BAF dataset, the frozen cost
matrix, the temporal split, and the delay regimes tested.

% =============================================================
\section{Conclusion and Future Work}
\label{sec:conclusion}
% =============================================================

We have presented a cost-sensitive fraud decision policy, evaluated it
end-to-end under a delayed and censored label regime on the Bank Account
Fraud dataset, and tested its operational robustness with a pre-registered
falsification study. The policy reduces realized cost per transaction by
$57.91\%$ relative to the strongest static baseline, with a $95\%$ bootstrap
confidence interval of $[53.95\%, 62.13\%]$ that excludes zero. The result
is robust to a $2\times$ variation in each cost parameter, with a minimum
advantage of $47.76\%$ across the sweep. It survives constrained review
capacity down to $2\%$ of test volume, survives a two-month label delay,
and is confirmed under rolling retraining with a fragility note. Two
proposed concerns --- group disparity and stratified calibration --- are
tested and killed as non-findings.

The falsification study is the paper's methodological contribution. It
shows that a headline result can be validated operationally without
retraining on the test set, without re-touching the frozen evaluation
window, and without expanding the modeling scope. Two of the six
hypotheses are reported as non-findings at the same length as the
survivors.

Several directions are named but not pursued. Capacity-aware ranking and
scheduling in wall-clock time; multi-regime delay with sub-month
granularity; a drift-correction model; causal estimation under
counterfactual policy selection; and amount-scaled false-positive cost are
all natural extensions. Each would require its own pre-registered kill
criterion before implementation, following the discipline established here.

% =============================================================
\begin{thebibliography}{00}
% =============================================================

\bibitem{jesus2022} S.~Jesus, J.~Pombal, D.~Alves, A.~Cruz, P.~Saleiro,
R.~Ribeiro, J.~Gama, and P.~Bizarro, ``Turning the tables: Biased,
imbalanced, dynamic tabular datasets for ML evaluation,'' in \emph{Proc.
NeurIPS Datasets and Benchmarks Track}, 2022.

\bibitem{elkan2001} C.~Elkan, ``The foundations of cost-sensitive
learning,'' in \emph{Proc. 17th Int. Joint Conf. Artificial Intelligence},
2001, pp.~973--978.

\bibitem{elkan2008} C.~Elkan and K.~Noto, ``Learning classifiers from only
positive and unlabeled data,'' in \emph{Proc. 14th ACM SIGKDD Int. Conf.
Knowledge Discovery and Data Mining}, 2008, pp.~213--220.

\bibitem{chapelle2014} O.~Chapelle, ``Modeling delayed feedback in display
advertising,'' in \emph{Proc. 20th ACM SIGKDD Int. Conf. Knowledge
Discovery and Data Mining}, 2014, pp.~1097--1105.

\bibitem{dalpozzolo2015} A.~Dal~Pozzolo, O.~Caelen, R.~A.~Johnson, and
G.~Bontempi, ``Calibrating probability with an unbalanced class: An
application to fraud detection,'' in \emph{Proc. IEEE Symp. Computational
Intelligence and Data Mining}, 2015, pp.~1--7.

\bibitem{pedregosa2011} F.~Pedregosa \emph{et al.}, ``Scikit-learn: Machine
learning in Python,'' \emph{Journal of Machine Learning Research}, vol.~12,
pp.~2825--2830, 2011.

\bibitem{ke2017} G.~Ke \emph{et al.}, ``LightGBM: A highly efficient
gradient boosting decision tree,'' in \emph{Proc. Advances in Neural
Information Processing Systems}, 2017, pp.~3146--3154.

\end{thebibliography}

% =============================================================
\appendices
\section{Reproducibility}
% =============================================================

Every result in this paper is reproducible from the raw dataset, the
configuration files, and a fixed seed ($42$). The pipeline executes in one
command. The code, the falsification plan, the pre-registered kill criteria,
the six intermediate evidence files, and the assembled falsification report
are released alongside the paper.

Test coverage is 68 unit tests spanning the policy, the cost matrix, the
data schema, the model, the app validation, and the two Phase~B
components. All pass on the released commit.

\end{document}